\section{Conclusiones}
Este capítulo tiene que incluir:
\begin{itemize}
\item Una descripción de las conclusiones del trabajo:
\begin{itemize}
    \item Una vez se han obtenido los resultados, ¿qué conclusiones se extraen?
    \item ¿Estos resultados son los esperados? ¿O han sido sorprendentes? ¿Por qué? 
\end{itemize}
\item Una reflexión crítica sobre el logro de los objetivos planteados inicialmente:
\begin{itemize}
    \item ¿Hemos logrado todos los objetivos? Si la respuesta es negativa, ¿por qué motivo?
\end{itemize}
\item Un análisis crítico del seguimiento de la planificación y metodología a lo largo del producto:
\begin{itemize}
    \item ¿Se ha seguido la planificación?
    \item ¿La metodología prevista ha sido suficientemente adecuada?
    \item ¿Ha habido que introducir cambios para garantizar el éxito del trabajo? ¿Por qué? 
\end{itemize}
\item De los impactos previstos en \ref{s:etic}, ético-sociales, de sostenibilidad y de diversidad, evaluar/mencionar si se han mitigado (si eran negativos) o si se han conseguido (si eran positivos). 
\item Si han aparecido impactos no previstos a \ref{s:etic}, evaluar/mencionar cómo se han mitigado (si eran negativos) o que han aportado (si eran positivos).
\item Las líneas de trabajo futuro que no se han podido explorar en este trabajo y han quedado pendientes.


%\item Una descripción de las conclusiones del trabajo: Qué lecciones se han aprendido del trabajo?
%\item Una reflexión crítica sobre el logro de los objetivos planteados inicialmente: Hemos conseguido todos los objetivos? Si la respuesta es negativa, por qué motivo?
%\item De los impactos previstos a la sección \ref{s:etic}, una evaluación, o al menos, mención, sobre si se han mitigado (si eran negativos) o si se han conseguido (si eran positivos). 
\end{itemize}

%\section{Líneas de futuro}
%Las líneas de trabajo futuro que no se han podido explorar en este trabajo y han quedado pendientes.

%\section{Seguimiento de la planificación}
%Un análisis crítico del seguimiento de la planificación y metodología a lo largo del trabajo: 
%\begin{itemize}
%    \item  Se ha seguido la planificación? 
%    \item La metodología prevista ha sido la adecuada? 
%    \item Ha habido que introducir cambios para garantizar el éxito del trabajo? ¿Por qué? 
%\end{itemize}

\section{Líneas futuras de trabajo}
Durante la fase de análisis de objetivos, se han tenido en cuenta las siguientes líneas de trabajo que pueden dar continuación y mejorar los hitos reflejados en el proyecto:
\\

\textbf{Línea 1. Construir un portal web informativo y de participación ciudadana}

Diseñar un portal web web con orientación pública para divulgar los planes y políticas de accesibilidad universal, habilitando un canal interactivo de consulta y participación ciudadana.
\\

\textbf{Línea 2. Incluir un catálogo de cumplimiento de accesibilidad por municipio}

Integrar en el portal ciudadano un catálogo de accesibilidad municipal que muestre información detallada de los planes de accesibilidad en marcha, así como el visor interactivo de cada localidad.
\\

\textbf{Línea 3. Habilitar la integración con estándares cartográficos de geoportales}

Estructurar los datos obtenidos tras la predicción como servicios web de mapas, permitiendo una integración directa con \gls{SIG} o geoportales.
\\

\textbf{Línea 4. Implementar un algoritmo de recomendación de plataformas únicas}

Diseñar un modelo prescriptivo que, a partir de las métricas extraídas, permita identificar vías inaccesibles que, según los criterios de la normativa legal vigente, sean candidatas para su transformación en plataformas únicas con prioridad peatonal.
\\

\textbf{Línea 5. Desarrollar un módulo abierto para la ingesta y procesamiento de imágenes}

Facilitar en el portal público un flujo para el proceso de imágenes aéreas o satelitales, habilitando a municipios pequeños y a la ciudadanía  la visualización y análisis de datos de sus localidades.
\\

\textbf{Línea 6. Incorporar en el modelo de accesibilidad factores temporales o dinámicos}

Existen factores temporales que, para acotar el alcance del proyecto no se han incorporado, relacionadas con el la característica dinámica de las ciudades. Entre otros factores encontrados se han encontrado: barreras y obstáculos en determinadas franjas horarias, como carga y descarga, aglomeración de personas en eventos, mercadillos, horario de depósito de cubos de basura en la calle, entrada y salida de colegios y trabajos, o la ocupación de aceras de los talleres mecánicos cuando están abiertos.
\\

\textbf{Línea 7. Mapa de rutas accesibles en las ciudades}

Modelar una red de caminos accesibles de un punto a otro de la ciudad, que permitan realizar un trayecto con un mínimo nivel de accesibilidad, permitiendo visualizar en un mapa la ruta, y la comparación con la ruta más corta que no cumple con los criterios, permitiendo medir la diferencia y sobreesfuerzo que se necesita para completar la misma ruta.
\\

\textbf{Línea 8. Sistema híbrido con imágenes a nivel de calle}

Realizar un sistema que contenga tanto modelos que analicen las imágenes aéreas, con imágenes a nivel de calle obtenidas a partir de softwares como Google Street View, o imágenes particulares, que permitan realizar incluir factores no visibles desde imágenes en dos dimensiones aéreas, como la altura de las aceras, la estructura de semáforos, o la calidad del pavimento.
\\